\documentclass[11pt]{beamer}
\usetheme{Madrid}
\usecolortheme{whale}

% Pakets
\usepackage[utf8]{inputenc}
\usepackage{graphicx}
\usepackage{booktabs}

% Title Page Information
\title[FuzzyFit Assistant]{FuzzyFit: Adaptive Workout Intensity \& Recovery Guide}
\subtitle{CENG 386 - Fuzzy Logic Final Project}
\author[Your Name]{Your Name}
\institute[University Name]{
  Department of Computer Engineering \\
  University Name
}
\date{June 12, 2026}

\begin{document}

% Slide 1: Title Page
\begin{frame}
  \titlepage
\end{frame}

% Slide 2: The Problem
\begin{frame}{The Problem with Traditional Systems}
  \begin{itemize}
    \item \textbf{Rigid Logic:} Traditional fitness schedules rely on strict calendars (e.g., "Monday is heavy leg day").
    \item \textbf{Human Ambiguity:} Our physical readiness is not black and white. It fluctuates based on sleep, stress, and nutrition.
    \item \textbf{The Solution:} We need a system that understands the "fuzzy" nature of human physiology to prevent overtraining and optimize recovery.
  \end{itemize}
\end{frame}

% Slide 3: System Architecture
\begin{frame}{System Architecture (Inputs \& Outputs)}
  The system utilizes the \textbf{Mamdani Inference Model}.
  
  \vspace{0.5cm}
  \textbf{4 Linguistic Inputs (Antecedents):}
  \begin{enumerate}
      \item Sleep Quality (0-10)
      \item Muscle Soreness (0-10)
      \item Energy Level (0-10)
      \item Stress Level (0-10)
  \end{enumerate}
  
  \vspace{0.3cm}
  \textbf{2 Actionable Outputs (Consequents):}
  \begin{itemize}
      \item Workout Intensity (0-100\%)
      \item Workout Volume (0-120 Min)
  \end{itemize}
\end{frame}

% Slide 4: Membership Functions
\begin{frame}{Fuzzification (Membership Functions)}
  We use a combination of Triangular (\texttt{trimf}) and Trapezoidal (\texttt{trapmf}) functions to map crisp inputs to fuzzy linguistic terms.
  
  \vspace{0.3cm}
  \textbf{Example: Sleep Quality}
  \begin{itemize}
      \item \textit{Poor:} [0, 0, 3, 5] (Trapezoidal)
      \item \textit{Average:} [4, 6, 8] (Triangular)
      \item \textit{Good:} [7, 9, 10, 10] (Trapezoidal)
  \end{itemize}
  
  \vspace{0.3cm}
  \textit{Note: This prevents sharp cutoffs, allowing a sleep value of 4.5 to be partially "Poor" and partially "Average".}
\end{frame}

% Slide 5: The Rule Base
\begin{frame}{The Rule Base (Knowledge Base)}
  The engine consists of \textbf{18 carefully designed rules} to eliminate "dead zones".
  
  \vspace{0.3cm}
  \begin{block}{Rule Example 1 (Extreme Fatigue)}
    \textbf{IF} Sleep is Poor \textbf{OR} Soreness is High \textbf{OR} Stress is High \\
    \textbf{THEN} Intensity is Very Light \textbf{AND} Volume is Low.
  \end{block}
  
  \begin{block}{Rule Example 2 (Optimal State)}
    \textbf{IF} Sleep is Good \textbf{AND} Soreness is Low \textbf{AND} Energy is Surplus \textbf{AND} Stress is Low \\
    \textbf{THEN} Intensity is Maximum \textbf{AND} Volume is High.
  \end{block}
\end{frame}

% Slide 6: Defuzzification & Control Surfaces
\begin{frame}{Defuzzification \& 3D Control Surfaces}
  \begin{itemize}
      \item \textbf{Method:} Centroid (Center of Gravity) is used to calculate the final crisp workout parameters.
      \item \textbf{3D Surfaces:} The system generates dynamic 3D control surfaces to visualize non-linear decision boundaries.
      \item For instance, visualizing how Intensity scales exponentially when Energy increases and Stress decreases.
  \end{itemize}
\end{frame}

% Slide 7: Technology Stack & UI/UX Justification
\begin{frame}{Technology Stack \& UI/UX (Python vs MATLAB)}
  While the mathematical foundation (Mamdani/Centroid) is identical to MATLAB's Fuzzy Logic Toolbox, \textbf{Python (scikit-fuzzy) + Streamlit} was chosen to create a modern web application.
  
  \vspace{0.3cm}
  \textbf{Gestalt UI/UX Principles Implemented:}
  \begin{itemize}
      \item \textbf{Law of Common Region:} Inputs and outputs are strictly encapsulated in visually bounded containers.
      \item \textbf{Proximity:} Technical engine settings are isolated in the sidebar.
      \item \textbf{Symmetry:} The layout balances data entry with data visualization equally.
  \end{itemize}
\end{frame}

% Slide 8: Live Demo & Q&A
\begin{frame}{Conclusion}
  \begin{center}
      \Large \textbf{Live Demo}
      
      \vspace{0.5cm}
      \normalsize
      Transitioning to the interactive web application to demonstrate the Fuzzy Inference System in real-time.
      
      \vspace{1cm}
      \Large \textbf{Thank You!} \\
      \normalsize Any Questions?
  \end{center}
\end{frame}

\end{document}
